\documentclass[11pt,leqno]{amsart}
\usepackage[letterpaper,margin=1in]{geometry}
\usepackage{amsmath,amssymb,mathtools}
\usepackage{microtype}
\usepackage{needspace}
\usepackage[colorlinks=true,linkcolor=blue,citecolor=blue,urlcolor=blue]{hyperref}
\hypersetup{pdftitle={The sharp deficit in primitive signed reciprocal relations},
 pdfauthor={Will Howe, Grant Huffman, Jensen Kohlmeyer, Liam Kruer},
 pdfsubject={Erdos Problem 319}}
\numberwithin{equation}{section}
\newtheorem{theorem}{Theorem}[section]
\newtheorem{proposition}[theorem]{Proposition}
\newtheorem{lemma}[theorem]{Lemma}
\newtheorem{corollary}[theorem]{Corollary}
\theoremstyle{definition}
\newtheorem{definition}[theorem]{Definition}
\theoremstyle{remark}
\newtheorem{remark}[theorem]{Remark}
\newcommand{\N}{\mathbb N}
\newcommand{\Z}{\mathbb Z}
\newcommand{\Q}{\mathbb Q}
\newcommand{\R}{\mathbb R}
\newcommand{\e}{\mathrm e}
\newcommand{\lcm}{\operatorname{lcm}}
\newcommand{\pp}{\mathcal Q}
\newcommand{\eps}{\varepsilon}
\newcommand{\ind}{\mathbf 1}
\newcommand{\Ph}{P^{+}}
\newcommand{\card}[1]{\lvert #1\rvert}
\newcommand{\norm}[1]{\left\lVert #1\right\rVert}
\newcommand{\rec}[1]{R(#1)}
\newcommand{\floor}[1]{\left\lfloor #1\right\rfloor}
\newcommand{\ceil}[1]{\left\lceil #1\right\rceil}
\newcommand{\dd}{\,\mathrm d}
\newcommand{\Logtwo}{\log\log}
\newcommand{\Logthree}{\log\log\log}
\allowdisplaybreaks[1]

\title[Primitive signed reciprocal relations]{The sharp deficit in primitive signed reciprocal relations}
\author{Will Howe}
\author{Grant Huffman}
\author{Jensen Kohlmeyer}
\author{Liam Kruer}
\thanks{Email addresses.\ 
\href{mailto:howew@purdue.edu}{\nolinkurl{howew@purdue.edu}} (Will Howe),
\href{mailto:huffma97@purdue.edu}{\nolinkurl{huffma97@purdue.edu}} (Grant Huffman).\newline
\href{mailto:jenwin@purdue.edu}{\nolinkurl{jenwin@purdue.edu}} (Jensen Kohlmeyer),
\href{mailto:kruerl@purdue.edu}{\nolinkurl{kruerl@purdue.edu}} (Liam Kruer).}
\date{November 12, 2026}
\subjclass[2020]{11D68, 11B75, 11N25}
\keywords{Egyptian fractions, primitive relations, reciprocal sums, Erd\H{o}s Problem 319}

\begin{document}
\begin{abstract}
Let $c(N)$ be the largest cardinality of a subset of $\{1,\ldots,N\}$
admitting a signed reciprocal sum equal to zero with no nonempty proper
zero subsum under the same signs. We prove
\[
 N-c(N)\sim \e^{-\gamma}
 \frac{N\log\log\log N}{\log\log N},
\]
where $\gamma$ is Euler's constant. The construction joins divisor blocks
along a path of primes and completes their reciprocal sum while avoiding
those primes. For the converse, we prove a sharp harmonic bound for sets
whose reciprocals have no nonempty zero subsum modulo an arbitrary
integer. Exact rational representation in a dense sign class then turns
this modular bound into a lower bound for the omitted denominators.
\end{abstract}
\maketitle

\section{Introduction}

For an integer $N\ge1$, write $[N]=\{1,\ldots,N\}$.
A pair $(A,\delta)$, where $A\subseteq[N]$ and
$\delta\colon A\longrightarrow\{-1,1\}$, is a \emph{primitive signed
reciprocal relation} if
\begin{equation}\label{eq:primitive}
 \sum_{n\in A}\frac{\delta(n)}n=0,
 \qquad
 \sum_{n\in B}\frac{\delta(n)}n\ne0
 \quad(\varnothing\ne B\subsetneq A).
\end{equation}
The signs in the second condition are the restrictions of $\delta$.
Let $c(N)$ be the maximum cardinality of such a support, with value zero
when no nonempty support exists. Erd\H{o}s and Graham asked for the size
of this extremal function \cite[p.~43]{EG}. The question is recorded as
Problem~319 in \cite{Bloom319}.

\begin{theorem}\label{thm:main}
As $N\longrightarrow\infty$ through the positive integers,
\begin{equation}\label{eq:main}
 N-c(N)=\bigl(\e^{-\gamma}+o(1)\bigr)
 \frac{N\log\log\log N}{\log\log N}.
\end{equation}
\end{theorem}

Throughout the paper, all logarithms are natural. We use the notation
\begin{equation}\label{eq:scales}
 L=\log N,\qquad \lambda=\log\log N,\qquad
 \mu=\log\log\log N,\qquad \Delta_N=\frac{N\mu}{\lambda}.
\end{equation}
These quantities are used only for sufficiently large $N$.
For a finite set $E$ of positive integers, put
\[
 R(E)=\sum_{n\in E}\frac1n,
 \qquad E_q=\{n\in E:q\mid n\},\qquad D_E=\lcm\{n:n\in E\},
\]
and set $D_\varnothing=1$. We write $\pp(E)$ for the set of all prime
powers dividing at least one member of $E$. In particular, every positive
power of a prime up to its exponent in $D_E$ belongs to $\pp(E)$.
An integer is \emph{$S$-powersmooth} if every prime-power divisor is at
most $S$. We take $\Ph(1)=1$ and use $\Omega(n)$ for the number of prime
factors counted with multiplicity.

Croot's short-interval theorem \cite{Croot} gives the earlier lower bound
$c(N)\ge(1-1/\e+o(1))N$, as observed by Adenwalla
\cite{Bloom319}. A density-one construction was subsequently proposed
by g-01234, also known as popular\_12345 \cite{DensityOne}. Its mechanism
couples a positive family with a dense set of denominators avoiding a
specified collection of primes. We use this mechanism with a path of
pairwise connectors and quantitative denominator control. A complementary
large-prime obstruction appears in \cite{RexHannes}. The exact rational
representation theorem of Liu and Sawhney \cite{LS} supplies the analytic
input for our converse. We state the form needed here and give its
divisor and Fourier argument in Appendix~\ref{app:representation}.

\section{Analytic preliminaries}\label{sec:prelim}

We use the prime number theorem, partial summation, and the Mertens
estimates in the forms
\begin{equation}\label{eq:mertens}
 \sum_{p\le x}\frac1p=\log\log x+O(1),\qquad
 \prod_{p\le x}(1-p^{-1})^{-1}\sim\e^\gamma\log x,
 \qquad \sum_{p^j\le x}\log p\sim x.
\end{equation}
Here and below a letter $p$ in a prime sum denotes a prime, and $j\ge1$.
The usual prime-number-theorem error term also gives
$\sum_{x<p\le x^2}p^{-1}=\log2+o(1)$.

\begin{lemma}\label{lem:cleanestimates}
Put $S=N\exp(-6L^{4/5})$. Then
\begin{align}
 \Theta(N,S):=\sum_{\substack{q\le N\\q\text{ a prime power}\\q>S}}
 \frac1q&=O(L^{-1/5}),\label{eq:pptail}\\
 \card{\{n\le N:n\text{ is not }S\text{-powersmooth}\}}
 &\le N\Theta(N,S),\label{eq:nonsmooth}\\
 \card{\{n\le N:\Omega(n)>5\lambda\}}&=O(NL^{-3/2}).
 \label{eq:omega}
\end{align}
If $Y\longrightarrow\infty$ and $u\ge1$, then
\begin{equation}\label{eq:rankin}
 \sum_{\substack{n>Y^u\\\Ph(n)\le Y}}\frac1n
 \ll \e^{-u}\log Y.
\end{equation}
The implied constants are absolute.
\end{lemma}

\begin{proof}
For primes in $(S,N]$, partial summation and $\pi(x)\ll x/\log x$ give
\[
 \sum_{S<p\le N}\frac1p
 \ll \frac{1+\log(N/S)}{\log S}\ll L^{-1/5}.
\]
For powers of exponent at least two, the sum over bases $p\le\sqrt S$
is $O(S^{-1/2})$, since each geometric tail is at most $2/S$.
The bases $p>\sqrt S$ contribute at most
$2\sum_{n>\sqrt S}n^{-2}=O(S^{-1/2})$.
This proves \eqref{eq:pptail}, and a union bound proves
\eqref{eq:nonsmooth}.

The multiplicative identity
\[
 (3/2)^{\Omega(n)}=\sum_{d\mid n}g(d),\qquad
 g(p^j)=\tfrac12(3/2)^{j-1}
\]
and its Euler product give
\[
 \sum_{n\le N}(3/2)^{\Omega(n)}
 \le N\prod_{p\le N}\left(1+\frac{1/2}{p-3/2}\right)
 \ll N L^{1/2}.
\]
Markov's inequality proves \eqref{eq:omega}, because
$1/2-5\log(3/2)<-3/2$.

For \eqref{eq:rankin}, take $\sigma=1/\log Y$. Rankin's inequality gives
\[
 \sum_{\substack{n>Y^u\\\Ph(n)\le Y}}\frac1n
 \le \e^{-u}\prod_{p\le Y}(1-p^{-1+\sigma})^{-1}.
\]
The product on the right is $O(\log Y)$. Indeed, its logarithm differs
from that of the Mertens product by $O(1)$, using
$p^\sigma-1\le\e\sigma\log p$ and
$\sum_{p\le Y}(\log p)/p\ll\log Y$. The terms of degree at least two in
the logarithmic expansions have bounded total.
\end{proof}

\begin{lemma}[Exact rational representation]\label{lem:exact}
For all sufficiently large $N$, let
$L^{-3/100}\le c\le1$ and let $E\subseteq[cN,N]\cap\N$ satisfy
\begin{gather}
 \card E\ge N/L^2,\qquad
 n\text{ is }S\text{-powersmooth and }\Omega(n)\le5\lambda
 \quad(n\in E),\label{eq:exacthyp1}\\
 qR(E_q)\ge L^{-1/25}\quad(q\in\pp(E)),
 \qquad S=N\exp(-6L^{4/5}).\label{eq:exacthyp2}
\end{gather}
If $r>0$ is rational, its reduced denominator divides $D_E$, and
\begin{equation}\label{eq:targetrange}
 (1+L^{-1})r\le R(E)\le Lr,
\end{equation}
then $R(F)=r$ for some $F\subseteq E$.
The threshold is uniform over $c,E,r$.
\end{lemma}

This is the logarithmic-parameter specialization of the reciprocal-sum
criterion of Liu and Sawhney \cite[Proposition~5.2]{LS}.
Appendix~\ref{app:representation} supplies the specialization and the
details of the divisor extraction used in its proof.

\begin{lemma}[Cleaning]\label{lem:cleaning}
For all sufficiently large $N$ and every $E\subseteq[N]$ there is
$F\subseteq E$ such that all members
of $F$ are $S$-powersmooth, $\Omega(n)\le5\lambda$ on $F$, and
$qR(F_q)\ge L^{-1/25}$ for every $q\in\pp(F)$. Moreover,
\begin{equation}\label{eq:cleanloss}
 \card{E\setminus F}
 \ll N\bigl(L^{-1/5}+L^{-3/2}+L^{-1/25}\lambda\bigr).
\end{equation}
After the first two restrictions, the further loss in reciprocal mass is
$O(L^{-1/25}\lambda)$.
\end{lemma}

\begin{proof}
Apply the first two restrictions using Lemma~\ref{lem:cleanestimates}.
Whenever a prime power $q$ divides a member of the current set and
$qR(E_q)<\eta:=L^{-1/25}$, remove all its multiples.
A processed prime power never reappears. The total reciprocal mass
removed is at most
\[
 \eta\sum_{\substack{q\le N\\q\text{ a prime power}}}\frac1q
 =O(\eta\lambda).
\]
The number of removed denominators is at most $N$ times this mass.
The finite process terminates with the required properties.
\end{proof}

\section{Reciprocal sums modulo an integer}\label{sec:modular}

For $q\ge2$, a finite set $V$ of positive integers coprime to $q$ is
\emph{inverse-zero-free modulo $q$} if
\[
 \sum_{n\in U}n^{-1}\ne0\pmod q
 \qquad(\varnothing\ne U\subseteq V).
\]
The inverses in this definition are taken in $\Z/q\Z$.
The partial sums along any ordering of $V$, including the initial zero,
are distinct. Consequently,
\begin{equation}\label{eq:length}
 \card V\le q-1.
\end{equation}

\begin{lemma}[Prime inverse sums]\label{lem:primecover}
There is an absolute constant $C$ with the following property.
For sufficiently large $q$, write $\ell=\log q$ and
$Y=\ell^3(1+\log\ell)^2$. If $Y\le X\le q$ and a set contains at least
\[
 k_X=\ceil{C\ell X^{2/3}(\log X)^{-2/3}}
\]
distinct primes in $[X,2X)$ coprime to $q$, its inverse subset sums cover
every residue class modulo $q$.
\end{lemma}

\begin{proof}
Use exactly $k=k_X$ primes. Fix $0<h<q$. For each selected prime $p$,
choose $u_p\in(-q/2,q/2]$ with $u_p\equiv hp^{-1}\pmod q$ and write
$pu_p=h+j_pq$. Put $J=\floor{k\log X/(32\ell)}$.
Uniformly in the stated range, $J\longrightarrow\infty$ and $J<q-1$.
Each nonzero integer $h+jq$ with $|j|\le J$ has magnitude less than $q^2$
and has at most $2\ell/\log X$ prime divisors in $[X,2X)$.
Thus fewer than $k/2$ selected primes have $|j_p|\le J$.
For all the others,
\[
 \frac{|u_p|}{q}\ge\frac{J-1}{2X}
 \gg\frac{k\log X}{\ell X}.
\]
The modulus of the $h$th Fourier coefficient of a uniformly random
inverse subset sum is
\[
 \prod_p\left|\cos\frac{\pi u_p}{q}\right|
 \le\exp\left(-c_0\frac{k^3(\log X)^2}{\ell^2X^2}\right)
 <q^{-2}
\]
when $C$ is sufficiently large. We used
$|\cos\pi t|\le\exp(-2t^2)$ for $|t|\le1/2$.
Fourier inversion now gives probability at least
$q^{-1}(1-(q-1)q^{-2})>0$ at each residue.
This is the prime-inverse Fourier argument of Croot
\cite[Lemma~1 of the preprint]{Croot}, with a varying prime interval.
\end{proof}

\begin{lemma}\label{lem:largeprimeweight}
Uniformly over inverse-zero-free $V$ modulo $q$,
\[
 \sum_{\substack{n\in V,\ n\le q\\\Ph(n)>Y}}\frac1n=O(1),
 \qquad Y=(\log q)^3(1+\log\log q)^2.
\]
\end{lemma}

\begin{proof}
Write $n=mp$, where $p=\Ph(n)$, and partition $p>Y$ into dyadic
intervals $[X,2X)$. For a fixed $m$, at most $k_X$ primes can occur.
Otherwise one can reserve a prime, express its negative inverse using
$k_X$ of the remaining primes by Lemma~\ref{lem:primecover}, and multiply
by $m^{-1}$ to obtain a forbidden zero subsum.
Since $\Ph(m)\le2X$, the contribution of this interval is at most
\[
 \frac{k_X}{X}\prod_{p\le2X}(1-p^{-1})^{-1}
 \ll (\log q)X^{-1/3}(\log X)^{1/3}.
\]
The geometric sum over the intervals is
$O((\log q)Y^{-1/3}(\log Y)^{1/3})=O(1)$.
\end{proof}

\begin{lemma}\label{lem:localdensity}
For every $\beta>0$, all sufficiently large $q$ have the following
property, uniformly for integers $t\ge\log q$ and inverse-zero-free $V$
modulo $q$,
\begin{equation}\label{eq:localdensity}
 \card{\{n\in V:t/2<n\le t,\ \Ph(n)>\log q\}}<\beta t.
\end{equation}
\end{lemma}

\begin{proof}
We may assume $0<\beta<1$. Suppose the set in \eqref{eq:localdensity}
has at least $\beta t$ members. Apply Lemma~\ref{lem:cleaning} with
ambient parameter $t$. The resulting set $W\subseteq(t/2,t]$ has
$\card W\ge\beta t/2$ for all sufficiently large $t$, and
\[
 \beta/2\le R(W)\le1.
\]
Put $D=D_W$ and $\ell=\log q$. Every member of $W$ has a prime factor
greater than $\ell$. If $\mathcal P$ is the collection of such primes
dividing $D$, a union bound gives $\sum_{p\in\mathcal P}p^{-1}\ge\beta/2$.
Therefore
\begin{equation}\label{eq:largelcm}
 \log D\ge\sum_{p\in\mathcal P}\log p
 \ge\ell\log\ell\sum_{p\in\mathcal P}\frac1p
 \ge\tfrac\beta2\ell\log\ell.
\end{equation}
In particular $q/D\longrightarrow0$, uniformly in $t$.
Choose a positive integer $x$ divisible by $q$ such that
$\beta/8\le x/D\le\beta/4$. The hypotheses of
Lemma~\ref{lem:exact}, with ambient parameter $t$ and $c=1/2$, hold for
$W$ and $r=x/D$. Its cardinality condition and target inequalities hold
once $t$ is sufficiently large in terms of $\beta$.
It gives a nonempty $U\subseteq W$ with $R(U)=x/D$.
Since $(D,q)=1$ and $q\mid x$, reduction modulo $q$ contradicts inverse
zero-freeness. Every threshold in this argument is uniform for
$t\ge\log q$, proving the assertion.
\end{proof}

\begin{theorem}[Sharp modular harmonic bound]\label{thm:harmonic}
Uniformly over all finite inverse-zero-free sets $V$ modulo $q$,
\begin{equation}\label{eq:modularsharp}
 R(V)\le\bigl(\e^\gamma+o(1)\bigr)\log\log q
 \qquad(q\longrightarrow\infty).
\end{equation}
The coefficient $\e^\gamma$ is sharp, even when $q$ is restricted to primes.
\end{theorem}

\begin{proof}
Write $\ell=\log q$ and $Y=\ell^3(1+\log\ell)^2$.
The terms exceeding $q$ have total at most $1$ by \eqref{eq:length}.
Lemma~\ref{lem:largeprimeweight} treats the terms at most $q$ with
$\Ph(n)>Y$. For fixed $A>1$, Lemma~\ref{lem:cleanestimates} gives
\[
 \sum_{\substack{n>Y^A\\\Ph(n)\le Y}}\frac1n
 \ll\e^{-A}\log Y.
\]
There are $O(A\log Y)$ dyadic intervals between $\ell$ and $Y^A$.
By Lemma~\ref{lem:localdensity}, each contributes at most $2\beta$ to
the reciprocal mass of terms with $\Ph(n)>\ell$, for sufficiently large
$q$ in terms of $\beta$. Thus this entire part has mass
\[
 O(1)+O(\e^{-A}\log Y)+O(\beta A\log Y).
\]
Divide by $\log\ell$. First let $q$ tend to infinity with $A,\beta$
fixed, then let $A$ tend to infinity and $\beta A$ tend to zero.
The contribution is $o(\log\ell)$, uniformly in $V$.
The remaining terms have total at most
\[
 \sum_{\Ph(n)\le\ell}\frac1n
 =\prod_{p\le\ell}(1-p^{-1})^{-1}
 =\bigl(\e^\gamma+o(1)\bigr)\log\ell.
\]

For sharpness, let $q$ be prime and put
$D=\lcm(1,\ldots,\floor{\tfrac14\log q})$.
The set of positive divisors of $D$ is coprime to $q$ for large $q$.
Every nonempty reciprocal subtotal is $a/D$ with
$0<a\le DR(\{d:d\mid D\})<q$, using
$\log D\sim\tfrac14\log q$ and the Mertens product.
It is consequently inverse-zero-free, and its reciprocal mass is
$(\e^\gamma+o(1))\log\log q$.
\end{proof}

\begin{corollary}\label{cor:uniformdensity}
For every $\eps>0$, all sufficiently large $N$ satisfy the following
uniform assertion. If integers $q,t$ satisfy $2\le q\le N$ and
$t\ge\exp(\lambda^3)$, and
$V\subseteq[t]$ is inverse-zero-free modulo $q$, then
$\card V\le\eps t$.
\end{corollary}

\begin{proof}
For $q$ above a sufficiently large fixed threshold,
Lemma~\ref{lem:localdensity} and a dyadic partition bound the number of
members with $\Ph(n)>\log q$ by $\eps t/2$.
All remaining members are $L$-smooth. Rankin's inequality, with
$\sigma=1/\log L=1/\lambda$, gives
\[
 \card{\{n\le t:\Ph(n)\le L\}}
 \le t^{1-\sigma}\prod_{p\le L}(1-p^{-1+\sigma})^{-1}
 \ll t\lambda\e^{-\lambda^2}=o(t).
\]
For bounded $q$, use \eqref{eq:length}.
\end{proof}

\Needspace{20\baselineskip}
\section{A dense primitive construction}\label{sec:construction}

For a positive integer $G$, write
\[
 \Z_{(G)}=\{a/b\in\Q:(b,G)=1\}.
\]
Membership is understood after reduction, or equivalently through any
representation with denominator coprime to $G$.

\begin{lemma}[A path of prime blocks]\label{lem:path}
Let $D\ge1$, put $H=\sum_{d\mid D}d^{-1}$ and $K=DH$, and let
$p_1,\ldots,p_k$ be distinct primes, where $k\ge2$,
$p_i>K+1$, and $p_i\nmid D$. Suppose
\[
 0<B\le p_i\le M,\qquad \rho=\sum_{i=1}^k\frac1{p_i}\le\frac12,
 \qquad G=\prod_{i=1}^k p_i.
\]
There is a finite set $E$ of positive integers such that
\begin{gather}
 (e,G)>1\quad(e\in E),\qquad R(E)\in\Z_{(G)},\label{eq:path1}\\
 R(U)\in\Z_{(G)},\ U\subseteq E
 \quad\Longrightarrow\quad U=\varnothing\text{ or }U=E,\label{eq:path2}\\
 \max E\le\max\{MD,4kM^4\},\qquad
 H\rho\le R(E)\le H\rho+\rho/B.\label{eq:path3}
\end{gather}
\end{lemma}

\begin{proof}
Place the block $\{p_i d:d\mid D\}$ at vertex $i$ of a path.
For a denominator divisible by $p_i$ exactly once, its residue at $p_i$
means the reduction of $p_i$ times its reciprocal modulo $p_i$.
Let $\mathcal S_i\subseteq\mathbb F_{p_i}$ be the set of block subsums
of these residues, and let $s_i$ be the full-block residue.
After multiplication by $D$, every such subsum is an integer between
$0$ and $K<p_i$. Hence residue zero occurs only for the empty block,
residue $s_i$ occurs only for the full block, and
\[
 \card{\mathcal S_i}\le K+1<p_i,
 \qquad \mathcal S_i=s_i-\mathcal S_i.
\]

At each internal vertex choose $a_i$ such that
$-a_i\notin\mathcal S_i$. Prescribe residue $a_i$ for the incoming
edge and $-s_i-a_i$ for the outgoing edge. Both are nonzero.
At an endpoint prescribe residue $-s_i$ for its unique edge.
For an edge joining $p_i$ and $p_{i+1}$, the Chinese remainder theorem
provides a unit residue class $u$ modulo $p_ip_{i+1}$ so that the
denominator $p_ip_{i+1}u$ has the two prescribed residues.

Among the first $4k$ positive representatives of this class, a different
marked prime $p$ excludes at most $4k/p+1$ representatives.
The total excluded is at most $4k\rho+k-2<4k$.
Choose an unexcluded representative $u_i<4kp_ip_{i+1}$ and add the edge
denominator $p_ip_{i+1}u_i$. Each edge denominator is divisible by exactly
its two incident marked primes, each to exponent one. The block and edge
denominators are therefore all distinct.

Let $E$ be their union. The residue of its sum is zero at every marked
prime, proving \eqref{eq:path1}. Suppose $R(U)\in\Z_{(G)}$ for a subset
$U\subseteq E$. At an internal vertex, selecting exactly the incoming
edge would require $-a_i\in\mathcal S_i$. Selecting exactly the outgoing
edge would require $s_i+a_i\in\mathcal S_i$, which is equivalent to the
same forbidden condition by complement symmetry. Thus the two edge
indicators agree at each internal vertex. All edges are selected together
or all are absent. In the latter case every block is empty, and in the
former every block is full. The same conclusion holds when $k=2$.
This proves \eqref{eq:path2}.

The block denominators are at most $MD$, and the edge denominators are
less than $4kM^4$. Finally,
\[
 0\le R(E)-H\rho
 =\sum_{i=1}^{k-1}\frac1{p_ip_{i+1}u_i}
 \le\frac1B\sum_{i=1}^{k-1}\frac1{p_i}\le\frac\rho B.
\]
\end{proof}

\Needspace{9\baselineskip}
\begin{lemma}[Uniform rational completion]\label{lem:completion}
For every fixed $\eta>0$ there is a constant $C_\eta$ such that the
following holds for all sufficiently large integers $m$.
Let $G\ge1$ and let $r\in[\eta,2]\cap\Z_{(G)}$. Suppose every
prime-power component of the reduced denominator of $r$ is at most
$m/(\log m)^{14}$, and every prime divisor of $G$ exceeds both $C_\eta$
and $(\log m)^{13}$. Then
\[
 r=R(C),\qquad C\subseteq[m],\qquad (c,G)=1\quad(c\in C).
\]
\end{lemma}

\begin{proof}
For a sufficiently large prime power $q$, put $X=(\log q)^{12}$ and
$k=\ceil{(\log q)^{10}}$. There are at least $k$ primes in $[X,2X)$
coprime to $q$, by the prime number theorem. Lemma~\ref{lem:primecover}
applies to any $k$ of them, since its required number is $o(k)$.
Consequently every residue modulo $q$ is a sum of inverses of a subset
of these primes. The corresponding denominators $qp$ have total
reciprocal mass at most
\begin{equation}\label{eq:completioncost}
 \frac{k}{qX}\ll\frac1{q(\log q)^2}.
\end{equation}
Choose $q_0\ge5$ so large that this assertion holds above $q_0$ and the
sum of the right side of \eqref{eq:completioncost} over all integers
$q>q_0$ is less than $\eta/2$.

Starting from $r$, choose the largest prime-power component $q>q_0$ of
its current reduced denominator, and write the current fraction as
$a/(qv)$, where $(q,v)=1$. Select distinct primes $p\in[X,2X)$ with
\[
 \sum_p p^{-1}\equiv av^{-1}\pmod q,
\]
and subtract $\sum_p1/(qp)$. The new denominator divides
$\lcm(v,\{p\})$. To see this, multiply the difference by
$q\lcm(v,\{p\})$ and use the congruence to cancel the full factor $q$.
All its prime-power components are strictly less than $q$.
The processed components thus form a strictly decreasing finite sequence.
The total subtracted mass is less than $\eta/2$, so the residual remains
positive.

Every chosen denominator has largest prime-power component exactly its
processed $q$. Thus different steps use disjoint denominators.
All introduced primes are at most $2(\log m)^{12}< (\log m)^{13}$, and
every chosen denominator is at most
$2m/(\log m)^2<m$ for large $m$. The denominator remains coprime to $G$
at every step.

At termination the residual has denominator dividing
$D_0=\lcm(1,\ldots,\floor{q_0})$ and lies in $(0,2]$.
For each of $2D_0$ successive disjoint intervals, Croot's theorem
\cite{Croot} supplies a finite set of distinct denominators with
reciprocal sum $1$. Scaling all these sets by $D_0$ gives disjoint
representations of $1/D_0$. Their denominators are fixed and divisible
by $30$. They represent every $j/D_0$ with $1\le j\le2D_0$ by taking
unions. This finite reservoir is disjoint from the previously used
denominators, which have at most two distinct prime factors.
Choose $C_\eta$ larger than every reservoir denominator. For large $m$
the reservoir lies in $[m]$ and is coprime to $G$, completing the proof.
\end{proof}

\begin{proposition}\label{prop:upper}
There are primitive signed reciprocal relations in $[N]$ with
\[
 N-\card A\le\bigl(\e^{-\gamma}+o(1)\bigr)\Delta_N.
\]
\end{proposition}

\begin{proof}
Put
\[
 s=\floor{L/1000},\quad D=\lcm(1,\ldots,s),\quad
 H=\sum_{d\mid D}\frac1d,\quad B=\floor{N^{1/20}}.
\]
The prime number theorem gives $D\le N^{1/100}$ for large $N$.
Moreover,
\begin{equation}\label{eq:divisormass}
 H=\prod_{p\le s}\frac{1-p^{-1-\floor{\log s/\log p}}}{1-p^{-1}}
 =\bigl(\e^\gamma+o(1)\bigr)\lambda.
\end{equation}
Indeed, the sum of the omitted numerator factors is at most
$\pi(s)/s=O(1/\log s)$, and the denominator is the Mertens product.

Choose the shortest initial segment $p_1,\ldots,p_k$ of the primes in
$(B,B^2]$ for which
\begin{equation}\label{eq:rhochoice}
 \frac{\log H+1}{H}\le\rho:=\sum_i\frac1{p_i}
 <\frac{\log H+1}{H}+\frac1B.
\end{equation}
The choice is possible because the reciprocal mass of the band tends
to $\log2$. Eventually $k\ge2$, $\rho\le1/2$, $DH+1<B$, and all the
chosen primes are coprime to $D$. Apply Lemma~\ref{lem:path} and put
$G=\prod_i p_i$. It gives a positive family $E$ with
\begin{equation}\label{eq:Eupper}
 \max E\le N^{3/4},\qquad R(E)=\log H+1+o(1).
\end{equation}
For the first inequality, the block bound is
$B^2D\le N^{11/100}$, and the edge bound follows from $k\le B^2$ and
$4k(B^2)^4\le4N^{1/2}$.

Let $n=\floor{N/H}$ and let $F$ consist of the $S$-powersmooth integers
in $(n,N]$ coprime to $G$, where $S=N\exp(-6L^{4/5})$.
By a union bound and Lemma~\ref{lem:cleanestimates},
\begin{equation}\label{eq:Fcount}
 N-\card F\le n+N\rho+O(NL^{-1/5}).
\end{equation}
We also have
\begin{equation}\label{eq:Fmass}
 R(F)=\log H+o(1).
\end{equation}
In fact, the unrestricted interval has reciprocal mass
$\log(N/n)+o(1)=\log H+o(1)$.
The powersmooth restriction removes mass at most
$O(NL^{-1/5}/n)=O(HL^{-1/5})=o(1)$.
For a marked prime $p\le B^2=o(n)$, the reciprocal mass of its multiples
in $(n,N]$ is at most
$p^{-1}(1+\log(2N/n))$. Summing this bound removes at most
$O(\rho(1+\log H))=o(1)$.

The difference $r=R(E)-R(F)$ tends to $1$ and belongs to $\Z_{(G)}$.
Its reduced denominator is $S$-powersmooth, because this holds for all
denominators in $E\cup F$ and is preserved by taking an lcm and reducing
a fraction. For large $N$,
\[
 S\le\frac{n}{(\log n)^{14}},\qquad
 \min_i p_i>\max\{C_{1/2},(\log n)^{13}\}.
\]
Lemma~\ref{lem:completion} gives $C\subseteq[n]$, coprime to $G$, with
$R(C)=r$. The sets $E,F,C$ are disjoint. Assign positive signs to $E$
and negative signs to $F\cup C$.

If a subfamily has signed sum zero, its positive reciprocal subtotal
belongs to $\Z_{(G)}$. Lemma~\ref{lem:path} forces its positive part
to be empty or all of $E$. Positivity of the unsigned summands then
forces its negative part to be empty or all of $F\cup C$, respectively.
The relation is primitive.

Finally, \eqref{eq:Fcount} and \eqref{eq:rhochoice} imply
\[
 N-\card{E\cup F\cup C}
 \le\frac{N}{H}\bigl(\log H+2+o(1)\bigr)+O(NL^{-1/5}).
\]
Now $H\sim\e^\gamma\lambda$ and $\log H\sim\mu$.
The last error is $o(\Delta_N)$, proving the proposition.
\end{proof}

\section{The lower bound for arbitrary primitive relations}\label{sec:lower}

\subsection{The smaller sign class}

\begin{lemma}\label{lem:minority}
Let $A=A_+\sqcup A_-$ be a primitive signed reciprocal relation in $[N]$
and put $d=N-\card A$. Uniformly over these relations,
\begin{equation}\label{eq:minority}
 \min\{\card{A_+},\card{A_-}\}
 \le (1+o(1))Nr+\frac{\sqrt r}{1-\sqrt r}\,d,
 \qquad r\sim200\frac\mu\lambda.
\end{equation}
In particular, if $d=O(\Delta_N)$, the smaller sign class has
$O(\Delta_N)$ elements.
\end{lemma}

\begin{proof}
Fix $b=1/100$ and put
\[
 \mathcal P=\{p:\lambda^2\le p\le L^b\},\qquad
 r=\prod_{p\in\mathcal P}(1-p^{-1}),\qquad \tau=L^{-b}.
\]
The Mertens product gives $r\sim(2/b)\mu/\lambda$.
We may suppose both sign classes have at least $Nr/2$ elements.
No $p\in\mathcal P$ can satisfy
\begin{equation}\label{eq:bothpopular}
 \card{(A_+)_p}\ge\tau N/p,
 \qquad \card{(A_-)_p}\ge\tau N/p.
\end{equation}
To prove this, restrict either sign class to $[NL^{-3b},N]$ and apply
Lemma~\ref{lem:cleaning}. The lower cutoff removes at most
$NL^{-3b}/p+1$ multiples of $p$.
The powersmooth and $\Omega$ restrictions remove
$O(NL^{-1/5}+NL^{-3/2})=o(NL^{-2b})$ integers.
Hence at least $\tau N/(2p)$ multiples of $p$ remain before pruning.
Their reciprocal mass is at least $\tau/(2p)\ge L^{-2b}/2$.
The pruning removes mass $O(L^{-4b}\lambda)=o(L^{-2b})$, so $p$ still
divides the lcm of the resulting set $U$.
Also $\card U\ge Nr/3$ and
\[
 r/3\le R(U)\le3b\lambda+O(1).
\]
Apply Lemma~\ref{lem:exact} to $U$ with target $1/p$.
Its mass inequalities follow from $p\ge\lambda^2$ and $p\le L^b$.
Doing this for both signs produces two subsets of equal mass $1/p$.
Their union is a nonempty proper zero subsum, since the full mass of
either sign class is at least $r/2>1/p$ for large $N$.
This contradiction proves the assertion.

Assign each prime in $\mathcal P$ to a sign with the greater number of
multiples. A term bearing a prime assigned to its opposite sign lies
in an exceptional set of size at most
\[
 E=\tau N\sum_{p\in\mathcal P}\frac1p=o(Nr).
\]
Write $r_+,r_-$ for the products of $1-1/p$ over the two assigned
collections, so $r_+r_-=r$. Relabel the signs so $r_-\le\sqrt r$.
The product $Q=\prod_{p\in\mathcal P}p$ is $N^{o(1)}$.
Periodicity modulo $Q$ counts the integers bearing a prime from each
collection as
\[
 N(1-r_+)(1-r_-)+O(Q)\le d+E+O(Q).
\]
The positive terms outside the exceptional set avoid all minus-assigned
primes. Their number is at most
\begin{align*}
 \card{A_+}
 &\le Nr+N(1-r_+)r_-+E+O(Q)\\
 &\le Nr+\frac{r_-}{1-r_-}(d+E+O(Q))+E+O(Q).
\end{align*}
Since $Q=o(Nr)$, this proves \eqref{eq:minority}.
\end{proof}

\subsection{Denominator excess and reciprocal mass}

Fix a positive integer $D$ and put $\nu_p=v_p(D)$ and
\begin{equation}\label{eq:missingpowers}
 s_p=p^{\nu_p+1},\qquad
 \mathcal S_D=\{s_p:s_p\le N\},\qquad
 \rho_D=\sum_{s\in\mathcal S_D}\frac1s.
\end{equation}
For $n\ge1$ define
\begin{equation}\label{eq:excess}
 b(n)=\prod_p p^{\max\{v_p(n)-\nu_p,0\}},\qquad
 Q_b=\prod_{p\mid b}p^{\nu_p+v_p(b)}.
\end{equation}
If $b(n)=b$, then $n=Q_bm$, where $m\mid D_b$ and $(m,b)=1$, with
$D_b=D/\prod_{p\mid b}p^{\nu_p}$.
For any subset of this group,
\begin{equation}\label{eq:compatibility}
 D\sum_n\frac1n=\frac1b\sum_m\frac{D_b}{m}.
\end{equation}
Thus its reciprocal sum has reduced denominator dividing $D$ exactly
when the inverse sum of its cofactors is zero modulo $b$.

\begin{lemma}\label{lem:weighted}
For every $\eps>0$ there are constants $C_\eps$ and $\delta_\eps>0$
such that, for sufficiently large $N$, the following holds uniformly in
$D$ and $V\subseteq[N]$. Suppose $V$ has no member dividing $D$, and
the cofactors in each group $b(n)=b>1$ are inverse-zero-free modulo $b$.
If $\rho_D\le\delta_\eps$, then
\begin{equation}\label{eq:weighted}
 R(V)\le C_\eps+(\e^\gamma+\eps)\lambda\rho_D.
\end{equation}
\end{lemma}

\begin{proof}
Let $h(b)$ be the supremum of reciprocal masses of inverse-zero-free
sets modulo $b$. Theorem~\ref{thm:harmonic} gives
$h(b)\le(\e^\gamma+\eps')\log\log b$ for sufficiently large $b$,
and \eqref{eq:length} gives $h(b)\le1+\log b$ for every $b\ge2$.
The contribution of group $b$ is at most $h(b)/Q_b$.

Fix a large constant $Y$. The sum over groups supported entirely on
primes at most $Y$ is bounded by
\[
 \sum_{\substack{b>1\\\Ph(b)\le Y}}\frac{1+\log b}{b}<\infty,
\]
uniformly in $D,N$. For all other groups, choose $Y$ so large that the
sharp estimate for $h(b)$ applies. Since $b\le N$ for a nonempty group,
their total mass is at most $(\e^\gamma+\eps')\lambda$ times the sum of
$Q_b^{-1}$ over groups containing a prime greater than $Y$.
For a prime with $s_p\le N$, put
\[
 a_p=\sum_{j\ge1}p^{-\nu_p-j}
 =\frac1{s_p(1-1/p)}.
\]
The required sum of weights is bounded by
\[
 \left(\sum_{p>Y}a_p\right)\prod_{s_p\le N}(1+a_p)
 \le\left(1+\frac1{Y-1}\right)\rho_D\exp(2\rho_D).
\]
First choose $\eps'$ small and $Y$ large, and then choose
$\delta_\eps$ small. This proves \eqref{eq:weighted}.
\end{proof}

\subsection{A dense sign class and the removal of small subsums}

Consider a sequence of primitive relations with $d=O(\Delta_N)$.
By Lemma~\ref{lem:minority}, orient the signs so that the negative
class $B$ has $N-O(\Delta_N)$ elements and the positive class $P$ has
$O(\Delta_N)$ elements. Their common reciprocal mass satisfies
\begin{equation}\label{eq:commonmass}
 R(P)=R(B)\ge H_N-H_{d+\card P}
 \ge\mu-\log\mu-O(1)=(1-o(1))\mu,
\end{equation}
where $H_j=\sum_{i=1}^j1/i$ and $H_0=0$.

Fix $0<a<\e^{-3}$. Apply Lemma~\ref{lem:cleaning} to $B\cap(aN,N]$,
and call the resulting set $F$. Put $D=D_F$ and $\rho=\rho_D$.
Then
\begin{gather}
 \card{(B\cap(aN,N])\setminus F}=o(\Delta_N),\qquad
 R(F)=\log(1/a)+o(1),\label{eq:coremass}\\
 \card{(aN,N]\setminus F}=o(N),\qquad \rho=o(1).
 \label{eq:corerho}
\end{gather}
Only the last assertion needs explanation. Every integer in $(aN,N]$
divisible by a member of $\mathcal S_D$ is absent from $F$.
Let $Z(n)$ count these divisors. They have pairwise coprime bases, and
$a<1/4$. Hence
\[
 \sum_{aN<n\le N}Z(n)\ge\frac{N\rho}{4},\qquad
 \sum_{aN<n\le N}Z(n)^2\le N(\rho+\rho^2).
\]
For the first inequality use
$\floor x-\floor{ax}\ge x/4$ for $x\ge1$.
Cauchy--Schwarz gives at least $N\rho/(16(1+\rho))$ integers in this
union. Its size is $o(N)$ by \eqref{eq:corerho}'s first assertion, so
$\rho=o(1)$.

Put $u_N=(\log(1/a)+1)/L$. For large $N$, Lemma~\ref{lem:exact} and
\eqref{eq:coremass} imply
\begin{equation}\label{eq:dense-targets}
 \begin{gathered}
 r\in\Q,\quad u_N\le r\le2,\quad \operatorname{den}(r)\mid D\\
 \Longrightarrow\quad r=R(U)\text{ for some }U\subseteq F.
 \end{gathered}
\end{equation}
The lower mass inequality uses $\log(1/a)>3$, and the upper one follows
from $Lr\ge\log(1/a)+1$.

In each positive group with a fixed $b>1$, successively remove
inclusion-minimal nonempty subsets whose reciprocal sum has denominator
dividing $D$. By \eqref{eq:compatibility} and the cyclic partial-sum
argument, each has at most $b$ elements and therefore mass at most
$b/Q_b\le1$. Remove members of the group $b=1$ as singletons.
Let $Z\subseteq P$ be the union of all removed blocks.
If $R(Z)\ge u_N$, a first collection of whole blocks crossing $u_N$
would have mass in $[u_N,u_N+1)$ and denominator dividing $D$.
By \eqref{eq:dense-targets}, it would have an equal-mass representation
in $B$. This would give a nonempty proper zero subsum, since
$R(P)>u_N+1$ by \eqref{eq:commonmass}.
Consequently
\begin{equation}\label{eq:peeled}
 R(Z)<u_N,\qquad \card Z<Nu_N=o(\Delta_N).
\end{equation}
The remaining positive groups are inverse-zero-free on their cofactors.
Lemma~\ref{lem:weighted} now gives, for every fixed $\eps>0$,
\[
 (1-o(1))\mu\le C_\eps+(\e^\gamma+\eps)\lambda\rho+u_N.
\]
Since $\mu\longrightarrow\infty$, we obtain
\begin{equation}\label{eq:rholower}
 \liminf_{N\to\infty}\frac{\lambda\rho}{\mu}\ge\e^{-\gamma}.
\end{equation}

\subsection{Counting omitted denominators}

Let $\mathcal S=\{s\in\mathcal S_D:s\le S\}$, with
$S=N\exp(-6L^{4/5})$. Let $T$ be the set of $S$-powersmooth integers
in $(aN,N]$ divisible by exactly one member of $\mathcal S$.
Then
\begin{equation}\label{eq:Tlower}
 \card T\ge(1-2a)N\rho-N\rho^2-4S-2N\Theta(N,S).
\end{equation}
Indeed, let $Z_0(n)$ count the divisors from $\mathcal S$.
The pointwise inequality
$\ind_{\{Z_0=1\}}\ge2Z_0-Z_0^2$, the first-moment bound
\[
 \sum_{aN<n\le N}Z_0(n)
 \ge(1-a)N\sum_{s\in\mathcal S}\frac1s-2\card{\mathcal S},
\]
and the second-moment bound
\[
 \sum_{aN<n\le N}Z_0(n)^2
 \le N\left(\sum_{s\in\mathcal S}\frac1s+
                 \left(\sum_{s\in\mathcal S}\frac1s\right)^2\right)
\]
prove the estimate before the powersmooth restriction.
Use $\card{\mathcal S}\le S$,
$\rho-\sum_{s\in\mathcal S}s^{-1}\le\Theta(N,S)$, and
\eqref{eq:nonsmooth} to obtain \eqref{eq:Tlower}.

A member of $T$ cannot divide $D$. Since it is $S$-powersmooth, every
first missing power that divides it belongs to $\mathcal S$.
Its excess therefore has the form $b=p^j$, and
$Q_b=p^{\nu_p+j}\le S$.
For a group of $P\setminus Z$ with this property, its cofactors lie in
$[\floor{N/Q_b}]$ and are inverse-zero-free modulo $b$.
Moreover,
\[
 \floor{N/Q_b}\ge\tfrac12\exp(6L^{4/5})
 \ge\exp(\lambda^3)
\]
for large $N$. Corollary~\ref{cor:uniformdensity} shows that for each
fixed $\eps>0$ its cardinality is at most $\eps N/Q_b$.
Summing over all such groups yields
\begin{equation}\label{eq:positiveT}
 \card{T\cap(P\setminus Z)}
 \le\eps N\sum_{s_p\le N}\sum_{j\ge1}p^{-\nu_p-j}
 \le2\eps N\rho.
\end{equation}

Every negative member of $T$ lies in $(B\cap(aN,N])\setminus F$.
Every other member is omitted from $A$ or lies in $Z$ or in
$T\cap(P\setminus Z)$. Combining
\eqref{eq:coremass}, \eqref{eq:peeled}, \eqref{eq:Tlower}, and
\eqref{eq:positiveT} gives
\begin{equation}\label{eq:finalcount}
 d\ge(1-2a-2\eps)N\rho-N\rho^2-o(\Delta_N).
\end{equation}
Here $S=o(\Delta_N)$ and $N\Theta(N,S)=o(\Delta_N)$ by
Lemma~\ref{lem:cleanestimates}.
Since $\rho=o(1)$, equations \eqref{eq:rholower} and
\eqref{eq:finalcount} imply
\[
 \liminf\frac d{\Delta_N}\ge(1-2a-2\eps)\e^{-\gamma}.
\]
Let $a$ and $\eps$ tend to zero after taking the lower limit in $N$.
This proves the lower bound for every sequence with $d=O(\Delta_N)$.
For extremizers this size restriction follows from
Proposition~\ref{prop:upper}. That proposition and the lower bound
complete the proof of Theorem~\ref{thm:main}.


\appendix
\section{Exact rational representation}\label{app:representation}

We prove Lemma~\ref{lem:exact} by the divisor extraction and Fourier
method of Liu and Sawhney \cite[Section~5]{LS}.
Throughout this appendix put
\begin{gather*}
 L=\log N,\quad \lambda=\log L,\quad z=L^{4/5},\quad
 \eta=L^{-1/25},\quad S=N\e^{-6z},\\
 M=\floor{cN/20000},\quad K=M\e^{-z},\quad w=\log(N/M),\\
 \Gamma=\max\left\{\frac\eta{L^{1/5}\lambda^3},
              \frac{\eta^2L^{3/5}}{w^2\lambda^5}\right\}.
\end{gather*}
Here $c\ge L^{-3/100}$, so $w=O(\lambda)$ uniformly. In particular,
\begin{equation}\label{eq:appendixparameters}
 \Gamma\gg\frac{L^{13/25}}{\lambda^7},\qquad
 \frac{\Gamma^2}{L^{1/50}(w+z)}
 \gg\frac{L^{11/50}}{\lambda^{14}}\longrightarrow\infty.
\end{equation}
Also $N^{.99}\le S\le K\le M\le N/10^4$ for all sufficiently large
$N$. All thresholds below are uniform over the stated range of $c$.

We use the following standard upper-bound sieve consequence, also in
\cite[Lemma~2.4]{LS}. If $\mathcal P$ is a collection of primes at most
$y$, and
$y\le\exp(\log X/\sqrt{\log\log X})$, then
\begin{equation}\label{eq:sieve}
 \card{\{n\in[X,2X)\cap\N:p\nmid n\ (p\in\mathcal P)\}}
 \ll X\prod_{p\in\mathcal P}(1-p^{-1}).
\end{equation}
The constant is absolute, and the same bound applies to a subinterval
by enlarging to the displayed interval.

\subsection{Divisor extraction}

\begin{lemma}\label{lem:divisorextraction}
Let $T\subseteq[M,N]\cap\N$ satisfy $\Omega(n)\le5\lambda$ for every
$n\in T$. Suppose $q\le K$ is a prime power and $qR(T_q)\ge\eta/2$.
Put
\[
 H_* =\exp\left(\frac{\eta z}{2\lambda^3w}\right).
\]
There are a positive integer $d$ and a nonempty $T^*\subseteq T_{qd}$
such that
\begin{equation}\label{eq:extraction}
 qd\ge K,\qquad \min T^*\ge H_*qd,\qquad
 qdR(T^*)\gg\frac\eta{L^{1/5}\lambda}.
\end{equation}
\end{lemma}

\begin{proof}
Set $y=\exp(z/(10\lambda))$. We have $2\le H_*\le y$ for large $N$.
For $n\in T_q$, define
\[
 d_n=\prod_{\substack{p\mid n/q\\p>y}}p^{v_p(n/q)}.
\]
Then $qd_n\mid n$ and
\[
 qd_n\ge n/y^{\Omega(n)}\ge M\e^{-z/2}\ge K.
\]
Call $n$ poor if it has fewer than two distinct prime divisors in
$[H_*,y]$. If $n$ is poor, so is $n/q$.

We count poor integers in $[M/q,N/q]$ using doubling intervals
$[X,2X)$. For integers with no prime in $[H_*,y]$,
\eqref{eq:sieve} and the Mertens estimates give
$O(X\log H_*/\log y)$.
For those with exactly one such prime $p$, divide by $p$ and apply the
sieve to the other primes of the interval. Leaving $p$ unsieved changes
the product by at most a factor $2$, and permits all its higher powers.
Summation over $p$ bounds their number by
$O(X\lambda\log H_*/\log y)$.
All sieve applications are valid, since $X\ge M/q\ge\e^z$ and
$\log y=z/(10\lambda)$, including after division by $p\le y$.

There are $O(w)$ doubling intervals, so their reciprocal mass gives
\[
 qR(\{n\in T_q:n\text{ poor}\})
 \ll w\lambda\frac{\log H_*}{\log y}
 \ll\frac\eta\lambda.
\]
Discard the poor elements. The remaining set has $q$ times reciprocal
mass at least $\eta/4$. Each of its members has two distinct primes in
$[H_*,y]$. Division by the prime power $q$ removes at most one of them,
and division by $d_n$ removes none. Consequently $n\ge H_*qd_n$.

Partition these members according to $d_n=d$. The possible $d$ have all
prime divisors in $(y,N]$, and
\[
 \sum_d\frac1d\le\prod_{y<p\le N}(1-p^{-1})^{-1}
 \ll\frac L{\log y}\ll L^{1/5}\lambda.
\]
Since $qR(T_q\setminus\{\text{poor elements}\})$
is the sum of $d^{-1}qdR(T^*)$ over the fibers, one fiber satisfies
\eqref{eq:extraction}.
\end{proof}

The quotient valuation $v_p(n/q)$ in this proof ensures $qd_n\mid n$.
The use of two distinct intermediate primes ensures that one survives
division by the prime power $q$. Both details are needed for
\eqref{eq:extraction}.

\subsection{Fourier inversion and small frequencies}

Let $E$ satisfy Lemma~\ref{lem:exact}, put $Q=D_E$, and set
$\tau=r/R(E)$. The target inequalities give
\[
 L^{-1}\le\tau\le\frac L{L+1},\qquad
 \tau(1-\tau)\ge\frac1{2L}.
\]
Independently select every $n\in E$ with probability $\tau$ and call
the resulting subset $U$. Write $\mathrm e(t)=\exp(2\pi it)$ and
\[
 \Phi(h)=\prod_{n\in E}(1-\tau+\tau\mathrm e(h/n)).
\]
Since $Qr$ is an integer, Fourier orthogonality gives
\begin{equation}\label{eq:Fourierprob}
 \mathbb P(R(U)-r\in\Z)
 =\frac1Q\sum_{-Q/2<h\le Q/2}
       \operatorname{Re}\bigl(\mathrm e(-hr)\Phi(h)\bigr).
\end{equation}

The small-frequency estimate of \cite[Lemma~3.1]{LS} gives
\begin{equation}\label{eq:major}
 \sum_{|h|\le M/2}\operatorname{Re}
       \bigl(\mathrm e(-hr)\Phi(h)\bigr)\ge\frac34.
\end{equation}
Its hypotheses hold because $M\ge N^{.95}$,
$\card E\ge N/L^2\ge N^{.95}$, and
$\tau\in[L^{-2},1-L^{-2}]$ for large $N$.
For completeness, the estimate follows by expanding at $h=0$ after the
linear term is canceled by $r=\tau R(E)$. The leading term is the
positive Gaussian
\[
 \exp\left(-2\pi^2\tau(1-\tau)h^2
                  \sum_{n\in E}n^{-2}\right).
\]
The cubic remainder is uniform for $|h|\le M^{3/5}$ and is negligible
on the Gaussian scale. For $M^{3/5}<|h|\le M/2$, the inequality
\[
 |1-\tau+\tau\mathrm e(t)|\le
 \exp(-8\tau(1-\tau)\norm{t}^2)
\]
bounds the total contribution by $o(1)$. Here $\norm{t}$ denotes
distance to the nearest integer. These estimates give
\eqref{eq:major} uniformly in the present parameter range.

\subsection{Localization of large frequencies}

For an integer $h$, let $h_n\in(-n/2,n/2]$ be its least absolute
residue modulo $n$. Put
\[
 T_0=\frac{50N^2L\lambda}{\tau(1-\tau)K^2},\qquad
 I_h=(h-K/2,h+K/2),
\]
and define
\[
 \mathcal D_h=\left\{q\in\pp(E):
       \card{\{n\in E_q:|h_n|\ge K/2\}}<T_0\right\}.
\]
Each $n$ lies in exactly $\Omega(n)$ prime-power divisibility classes.
The preceding bound for a Fourier factor therefore gives
\begin{equation}\label{eq:Fourierdecay}
 |\Phi(h)|\le N^{-10\card{\pp(E)\setminus\mathcal D_h}}.
\end{equation}
Indeed, the product over all pairs $(q,n)$ with $q\mid n$ is at least
$|\Phi(h)|^{5\lambda}$. For every $q$ outside $\mathcal D_h$, at least
$T_0$ factors have $|h_n|\ge K/2$, and their product is at most
$\exp(-2\tau(1-\tau)K^2T_0/N^2)=\exp(-100L\lambda)$.
Taking the $1/(5\lambda)$ power proves an even stronger version of
\eqref{eq:Fourierdecay}.

We claim that there is an integer in $I_h$ divisible by every member
of $\mathcal D_h$. Fix $q\in\mathcal D_h$ and let
$T_q=\{n\in E_q:|h_n|<K/2\}$.
The definitions and the parameter bounds give
\[
 R(T_q)\ge\eta/q-T_0/M\ge\eta/(2q).
\]
For the last inequality it suffices to note that
\[
 \frac{S T_0}{M\eta}
 \ll c^{-3}L^2\lambda\eta^{-1}\e^{-4z}\longrightarrow0.
\]
Apply Lemma~\ref{lem:divisorextraction}. It produces $d_q$ and
$T_q^*\subseteq T_q$ with $qd_q\ge K$.
Every $n\in T_q^*$ has a multiple in $I_h$, which must be the unique
multiple of $qd_q$ in this interval. Call that integer $x_q$.

Let $E_q^*=\{n/(qd_q):n\in T_q^*\}$. Then
\begin{equation}\label{eq:scaledmass}
 \min E_q^*\ge H_*,\qquad
 \frac{\max E_q^*}{\min E_q^*}\le\e^w,\qquad
 R(E_q^*)\gg\frac\eta{L^{1/5}\lambda}.
\end{equation}
Partition $E_q^*$ into dyadic intervals $[2^j,2^{j+1})$, including the
one containing its minimum. If $\sigma=\log H_*$, the reciprocal
weights of these indices satisfy
\[
 \sum_j\frac1{j+1}\ll\min\{\lambda,w/\sigma\}.
\]
Weighted pigeonholing in \eqref{eq:scaledmass} gives a dyadic interval
$[y_q,2y_q)$ with
\begin{equation}\label{eq:dyadicmass}
 R(E_q^*\cap[y_q,2y_q))\ge\frac\Gamma{\log y_q}.
\end{equation}
To check the constants, the lower bound obtained before absorbing them
is an absolute positive multiple of
\[
 \frac1{\log y_q}
 \max\left\{\frac\eta{L^{1/5}\lambda^2},
             \frac{\eta^2L^{3/5}}{w^2\lambda^4}\right\}
 =\frac{\lambda\Gamma}{\log y_q}.
\]
Thus \eqref{eq:dyadicmass} follows for large $N$.
It also implies $\log y_q\gg\Gamma$, since the reciprocal mass of a
dyadic interval is bounded. Finally $y_q\le N/(qd_q)$.

For $q_1,q_2\in\mathcal D_h$, set $y=\min(y_{q_1},y_{q_2})$ and take
the prime interval
\[
 \mathcal P=\{p:\exp(L^{1/100})\le p
          \le\exp((\log y)L^{-1/100})\}.
\]
It is nonempty for large $N$ by \eqref{eq:appendixparameters}.
Let $\mathcal P_q$ consist of the primes in $\mathcal P$ dividing none
of the integers in $E_q^*\cap[y_q,2y_q)$.
The sieve \eqref{eq:sieve}, applied to the omitted primes, and
\eqref{eq:dyadicmass} imply
\[
 \sum_{p\in\mathcal P_q}\frac1p
 \le\log\left(\frac{C_0\log y_q}{\Gamma}\right)
\]
for an absolute $C_0$. The sieve cutoff holds because
$L^{-1/100}\le1/\sqrt{\log\log y_q}$ for large $N$.
The Mertens prime estimate gives
\begin{align*}
 \sum_{p\in\mathcal P\setminus(\mathcal P_{q_1}\cup\mathcal P_{q_2})}
 \frac1p
 &\ge\log\left(
 \frac{\Gamma^2}{2C_0^2 L^{1/50}
                         \log\max(y_{q_1},y_{q_2})}\right)\\
 &\ge\log\left(\frac{\Gamma^2}{2C_0^2 L^{1/50}(w+z)}\right)
 \ge1.
\end{align*}
Every prime counted here divides both $x_{q_1}$ and $x_{q_2}$.
Indeed it divides a cofactor $n/(qd_q)$, and $x_q$ is a multiple of $n$.
There are at least $\exp(L^{1/100})$ such primes, all at least that
large. Their product exceeds $N$, whereas
$|x_{q_1}-x_{q_2}|<K\le N$. Hence $x_{q_1}=x_{q_2}$.
This proves the claim. An empty $\mathcal D_h$ imposes no divisibility
condition.

\subsection{Summing the large frequencies}

For $M/2<|h|\le Q/2$ the set $\mathcal D_h$ is a proper subset of
$\pp(E)$. Otherwise its common multiple in $I_h$ would be a multiple
of $Q$, at distance less than $K/2$ from such a centered $h$, which is
impossible.
For fixed $\mathcal D\subsetneq\pp(E)$, the number of residue classes
$h$ for which $I_h$ contains a multiple of $\lcm(\mathcal D)$ is at most
\[
 (K+1)\frac Q{\lcm(\mathcal D)}
 \le N\prod_{q\in\pp(E)\setminus\mathcal D}q
 \le N^{1+\card{\pp(E)\setminus\mathcal D}}.
\]
Using \eqref{eq:Fourierdecay} and at most $N^s$ choices of a complement
of size $s$ gives
\begin{equation}\label{eq:minor}
 \frac1Q\sum_{M/2<|h|\le Q/2}|\Phi(h)|
 \le\frac1Q\sum_{s\ge1}N^{1+s}N^sN^{-10s}
 \le\frac2{QN}.
\end{equation}
Equations \eqref{eq:major} and \eqref{eq:minor} show that
\eqref{eq:Fourierprob} is at least $1/(2Q)$.

Finally, $\mathbb E R(U)=r$ and each selected summand is at most $1/M$.
The elementary bounded-summand concentration inequality gives
\[
 \mathbb P(|R(U)-r|\ge1)\le2\exp(-M^2/(2N)).
\]
Since every prime-power component of $Q$ is at most $S$, the prime
number theorem gives $\log Q=O(S)$. Also $S=o(M^2/N)$.
The last probability is therefore less than $1/(4Q)$ for large $N$.
With positive probability $R(U)-r$ is an integer of absolute value less
than $1$, so it is zero. This proves Lemma~\ref{lem:exact}.

\section*{Acknowledgments}
The authors acknowledge assistance from OpenAI's ChatGPT with proof
development, checking the mathematical arguments, and preparing the
exposition.

\begin{thebibliography}{9}

\bibitem{Bloom319}
T.~F. Bloom, \emph{Erd\H{o}s Problem \#319},
\url{https://www.erdosproblems.com/319}, accessed November 11, 2026.

\bibitem{Croot}
E.~S. Croot III, \emph{On unit fractions with denominators in short
intervals}, Acta Arith. \textbf{99} (2001), no.~2, 99--114.
\href{https://doi.org/10.4064/aa99-2-1}{doi:10.4064/aa99-2-1}.
Preprint \href{https://arxiv.org/abs/math/9904181}{arXiv:math/9904181}.

\bibitem{EG}
P.~Erd\H{o}s and R.~L. Graham, \emph{Old and New Problems and Results in
Combinatorial Number Theory}, Monographies de L'Enseignement
Math\'ematique, vol.~28, Universit\'e de Gen\`eve, Geneva, 1980.

\bibitem{DensityOne}
g-01234 (popular\_12345), \emph{A candidate proof that $c(N)/N\to1$ in
Erd\H{o}s Problem 319}, circulation draft, July 16, 2026.
\href{https://github.com/g-01234/erdos-319-density-one/blob/f55a25a29dcb17193918d27506be5e7a5b5d7d40/docs/proof.md}
{Author's manuscript}.

\bibitem{LS}
Y.~P. Liu and M.~Sawhney, \emph{On further questions regarding unit
fractions}, Int. Math. Res. Not. IMRN \textbf{2026} (2026), no.~2,
rnaf382. \href{https://doi.org/10.1093/imrn/rnaf382}{doi:10.1093/imrn/rnaf382}.
Preprint \href{https://arxiv.org/abs/2404.07113}{arXiv:2404.07113}.

\bibitem{RexHannes}
RexHannes, \emph{A large-prime obstruction for Erd\H{o}s Problem \#319},
unpublished mathematical note, 2026.

\end{thebibliography}
\end{document}

